% ====================================================================
% ====================================================================
% ====================================================================
%  Frame Gantt (échelle en MOIS) — Calendrier prévisionnel de la thèse
%  PRÉREQUIS dans le préambule :  \usepackage{pgfgantt}
%
%  Chaque case = 1 mois. L'axe démarre en NOVEMBRE 2025 (mois 1) et
%  couvre 36 mois (nov 2025 -> oct 2028). La 1re ligne = année de thèse
%  (nov->oct), la 2e = initiale du mois (N D J F M A M J J A S O).
%  Le trait pointillé orange = position actuelle (~ sept 2026, mois 11).
%  NB : les nombres des \ganttbar sont des indices de mois, 1 = nov 2025.
% ====================================================================

\section{Gantt : calendrier prévisionnel}
\begin{frame}[t]{Calendrier prévisionnel}

  \vspace{-2em}
  \begin{pccard}
    % Cale le graphe à gauche (au lieu de le centrer) : la marge négative
    % remonte les intitulés vers le bord gauche. Ajuste -1.2cm si besoin.
    \noindent\hspace*{-0.55cm}%
        \begin{ganttchart}[
        x unit=0.7cm,
        hgrid, vgrid={*3{draw=gray!12}, *1{draw=gray!25}},
        canvas/.style={draw=gray!40, line width=.5pt},
        y unit title=0.38cm, y unit chart=0.38cm,
        title height=1, title/.style={fill=none},
        title label font=\tiny\bfseries\color{pcNavy},
        group label font=\scriptsize\bfseries\color{pcNavy},
        bar label font=\scriptsize\color{pcNavy!80},
        group label node/.append style={left=0.15cm},
        bar label node/.append style={left=0.4cm},
        group/.append style={draw=none, fill=pcNavy!70},
        group height=.5, group peaks tip position=0,
        bar/.append style={draw=none, fill=pcBlue!75}, bar height=.6,
      ]{1}{12}

      % --- Bandeau 1 : année de thèse ---
      \gantttitle{Année 1 (2025--26)}{4}
      \gantttitle{Année 2 (2026--27)}{4}
      \gantttitle{Année 3 (2027--28)}{4}                       \\
      % --- Bandeau 2 : trimestres ---
      \gantttitle{T1}{1}\gantttitle{T2}{1}\gantttitle{T3}{1}\gantttitle{T4}{1}%
      \gantttitle{T1}{1}\gantttitle{T2}{1}\gantttitle{T3}{1}\gantttitle{T4}{1}%
      \gantttitle{T1}{1}\gantttitle{T2}{1}\gantttitle{T3}{1}\gantttitle{T4}{1} \\

      % ---------- Brique photovoltaïque ----------
      \ganttgroup{Semi-conducteur}{1}{7}                       \\
      \ganttbar{Solveur 1D (Van Roosbroeck)}{1}{7}             \\
      \ganttbar{Validation (Farrell, SCAPS)}{3}{7}             \\
      \ganttbar[bar/.append style={fill=pcSun!70}]%
               {Hétérojonctions \& CL de Robin}{5}{6}          \\
      \ganttbar[bar/.append style={fill=pcSun!70}]%
               {Validation données ANR}{7}{9}                  \\

      % ---------- Brique électrochimique ----------
      \ganttgroup{Électrochimie}{2}{9}                         \\
      \ganttbar{Modèle 3D $\rightarrow$ 1D (homogénéisation)}{2}{4} \\
      \ganttbar[bar/.append style={fill=pcSun!70}]%
               {Implémentation + modèle GDE}{5}{8}             \\
      \ganttbar[bar/.append style={fill=pcSun!70}]%
               {Validation données ANR}{8}{9}                  \\

      % ---------- Nouvelles briques ----------
      \ganttgroup{Nouvelles briques}{5}{8}                     \\
      \ganttbar[bar/.append style={fill=pcSun!70}]%
               {Photogénération}{5}{8}  \\
      \ganttbar[bar/.append style={fill=pcSun!70}]%
               {\textit{Detailed balance}}{7}{9}               \\

      % ---------- Couplage & optimisation ----------
      \ganttgroup{Couplage \& optimisation}{9}{12}             \\
      \ganttbar[bar/.append style={fill=pcSun!70}]%
               {Couplage des briques}{9}{11}                   \\
      \ganttbar[bar/.append style={fill=pcSun!70}]%
               {Outil d'optimisation}{10}{12}                  \\

      % ---------- Rédaction ----------
      \ganttbar[bar/.append style={fill=pcNavy!55}]%
               {Rédaction \& soutenance}{10}{12}               \\

      % ---------- Ligne "aujourd'hui" (~ sept. 2026, milieu du T4 année 1) ----------
      \ganttvrule[vrule/.style={pcSunDeep, line width=1pt, dashed},
                  vrule offset=0.5]%
                 {}{4}

    \end{ganttchart}

    \vspace{0.2em}
    \centering
    {\scriptsize\bfseries\color{pcNavy}
     Planning prévisionnel sur les trois années de thèse — %
     \textcolor{pcBlue!75}{$\blacksquare$}~réalisé \quad
     \textcolor{pcSun!70!black}{$\blacksquare$}~à venir \quad
     \textcolor{pcSunDeep}{$\dashv$}~aujourd'hui}
  \end{pccard}

\end{frame}